% Generated draft: requires mathematical and source review.
\documentclass[11pt]{article}
\usepackage[margin=25mm]{geometry}
\usepackage{fontspec}
\setmainfont{Latin Modern Roman}
\usepackage{amsmath,amssymb,mathtools,unicode-math}
\setmathfont{Latin Modern Math}
\usepackage{enumitem,needspace,xurl,hyperref}
\setlength{\parindent}{0pt}
\setlength{\parskip}{5pt}
\setlength{\emergencystretch}{3em}
\newcommand{\N}{\mathbb{N}}
\newcommand{\Z}{\mathbb{Z}}
\newcommand{\Q}{\mathbb{Q}}
\newcommand{\R}{\mathbb{R}}
\newcommand{\C}{\mathbb{C}}
\newcommand{\F}{\mathbb{F}}
\newcommand{\A}{\mathbb{A}}
\newcommand{\PP}{\mathbb{P}}
\newcommand{\E}{\mathbb{E}}
\newcommand{\eps}{\varepsilon}
\newcommand{\dd}{\,\mathrm{d}}
\begin{document}
\section*{UnsolvedMath ID: 30003495}
\textbf{Source problem code:} catalogue entry
\subsection{Definitions and mathematical statement}
Let \(V=V_1\otimes\cdots\otimes V_m\) be a tensor product of finite-dimensional modules over a quantum affine algebra, and let \(W\subseteq V\) be a common generalized eigenspace for the transfer matrices under consideration. Here a common generalized eigenspace means a subspace associated with fixed generalized eigenvalues of the operators in the intended transfer-matrix family. The source denotes by \(Q_{i,k}(z)\) the Baxter polynomials appearing in formulas for the joint eigenvalues associated with \(W\), with \(z\) their polynomial variable.

\textbf{Requested characterization.} Characterize the families of Baxter polynomials \(Q_{i,k}(z)\) that occur in these joint transfer-matrix eigenvalue formulas for \(W\).

\textbf{Unresolved formulation.} This is the mathematical target stated in the supplied record, but it is not a fully specified classification problem. The record does not specify the quantum affine algebra, its parameters, the coefficient field, the permitted modules, the transfer-matrix family or its eigenvalue formulas, the ranges and meanings of \(i\) and \(k\), or the definition and normalization of \(Q_{i,k}(z)\). Consequently, the polynomial families to be characterized cannot be defined precisely from the supplied evidence. No particular QQ-system, recurrence, degree restriction, root condition, or normalization is imposed here.
\subsection{Short English statement}
The problem asks which Baxter polynomials \(Q_{i,k}(z)\) appear in transfer-matrix eigenvalue formulas on a common generalized eigenspace in a tensor product of finite-dimensional quantum affine algebra modules. The supplied record leaves the polynomial family and the relevant formulas undefined, so a precise characterization problem requires further source clarification. Its open status is a supplied curation label.
\subsection{Sources}
\begin{itemize}
\item Local UnsolvedMath registry, permanent ID \(30003495\), source problem code \texttt{catalogue entry}, title \emph{Transfer Matrix Eigenvalue Polynomials for Quantum Affine Algebras}, in the Oberwolfach Reports Open Problems collection. No proposer is identified; the supplied proposed year is 2017.
\item Supplied report citation: \emph{Algebraic Groups (2018)}, Oberwolfach Reports, DOI \texttt{10.4171/owr/2017/21}, pp.\ 1281--1347. \url{https://doi.org/10.4171/owr/2017/21}. The background describes Oberwolfach Reports as documenting mathematical workshops and their open-problem discussions. The literature-source label for this same link is \emph{Oberwolfach Reports 2018}.
\item Supplied literature reference: \emph{Frenkel et al., q-Opers and QQ-Systems}. \url{https://arxiv.org/abs/2002.07344}.
\item Supplied literature reference: \emph{On Wronskians and QQ-Systems}. \url{https://arxiv.org/abs/2208.08018}.
\item The literature assessment dated \texttt{2026-08-21} is supplied dated curation, not independently verified here. It labels the problem open and reports that later papers construct many polynomial QQ-system solutions and Wronskian recurrences, while the extracted question does not precisely define the family \(Q_{i,k}\), preventing assignment of a general answer. The supplied background expressly calls for independent verification of time-sensitive status and cited claims.
\item The registry labels the statement \texttt{corrected\_verified} and reports that the official report was checked directly and that the clean statement restored definitions while preserving the target. These are supplied verification claims; this draft does not independently verify them or claim to have visited any source.
\end{itemize}
\end{document}
